\abstract

\keywords{BACKEND-ПЛАТФОРМА, МИКРОСЕРВИСНАЯ АРХИТЕКТУРА, ТРЕНЕР, СПОРТСМЕН, ТРЕНИРОВОЧНЫЙ ПРОЦЕСС, REST API, NATS JETSTREAM, POSTGRESQL, АНАЛИТИКА, AI-РЕКОМЕНДАЦИИ}

Выпускная квалификационная работа посвящена разработке backend-части открытой микросервисной платформы CoachLink для автоматизации тренировочного процесса и взаимодействия в системе <<тренер-спортсмен>>. В работе рассматривается проблема разрозненного обмена тренировочными заданиями и отчётами через мессенджеры, электронные таблицы и произвольные форматы данных.

Объектом исследования является процесс информационного взаимодействия между тренером и спортсменом при планировании, выполнении и анализе тренировок. Предметом исследования является backend-платформа для автоматизации назначения тренировочных заданий, сбора стандартизированных отчётов, отправки уведомлений и анализа данных тренировочного процесса.

Целью работы является разработка backend-части открытой микросервисной платформы для автоматизации тренировочного процесса и взаимодействия в системе <<тренер-спортсмен>>.

В рамках работы спроектирована и реализована микросервисная архитектура из семи независимых сервисов. Сервисы реализованы на языке Go, используют PostgreSQL для хранения данных, NATS JetStream для событийного взаимодействия и OpenAPI для документирования публичных интерфейсов. Backend обеспечивает регистрацию пользователей, разграничение ролей тренера и спортсмена, управление связями и тренировочными группами, назначение тренировок, приём стандартизированных отчётов, формирование уведомлений, агрегацию тренировочных данных и генерацию AI-рекомендаций на основе истории тренировок.

Первая версия ориентирована на лёгкую атлетику, однако архитектура позволяет адаптировать тренировки, отчёты и шаблоны под другие виды спорта.
